\chapter{本文工作及当前进展}
\section{总体研究框架与递进假设}
本文以“线索获取—对象绑定—意图确认—会话维持—动作反馈”为统一框架，研究三个逐层深入、但不互相替代的假设。第一，人的手持指向携带可用于选择物理目标的空间信息，方向性光学链路可在同一动作中完成目标发现与局部数据交换；RetroAR 已对该假设提供系统证据。第二，交互入口从手机迁移至眼镜后，用户能在较少手部操作下使用类似观察的动作选中目标，完成原位双向交互；EchoSight 已在相应任务中验证其效率与体验。第三，进一步测量眼动、主动对齐光束，并引入任务相关脑电，可能降低错选与误触发；分层链路可能改善持续会话，但这一假设须由 Omnisight 的完整任务实验检验。

三项工作采用同一分析顺序：\textbf{利用哪种注意力线索、如何绑定物理对象、如何通信与反馈、有哪些实验依据、留下什么问题}。RetroAR 和 EchoSight 的既有结果分别支持前两项假设，不直接证明第三项在神经信号、延迟或连接性上的收益。尤其不能把“眼镜朝向”混同于“实测眼动”，也不能把“看向”混同于“想操作”。

\section{RetroAR：从空间指向到目标绑定（已完成）}
\subsection{注意力线索与系统设计}
RetroAR 面向移动增强现实中的非接触定向交互\parencite{retroar}。用户举起手机，闪光灯照向目标端的逆反射标签 ViTag；标签用液晶矩阵调制反射光，手机后置摄像头接收返回信息。光束方向与摄像头视场共同约束候选对象，把“指向目标”和“发现目标”结合起来，无须先从无线列表查找对象。需要准确指出：系统利用的是手机空间指向这一\textbf{间接注意力代理}，没有测量用户实际视线。

为了让指向真正成为可用的交互入口，系统不仅要检测标签，还要在手机运动、视角变化和背景干扰下保持读取。动态兴趣区减少重复搜索的计算开销，几何检测帮助定位标签，自适应时空混合解码恢复信息，并协同支持六自由度位姿估计。这些设计共同服务于“选得中、读得出、跟得稳”，而非彼此独立的算法堆叠。

\subsection{已有证据与通向下一阶段的局限}
原论文报告最远约 $4\,\mathrm{m}$、最大约 $100^\circ$ 的工作范围，六自由度位姿估计的平移误差低至约 $1\,\mathrm{cm}$、旋转误差低至约 $4.7^\circ$；在所设非接触控制任务中，交互时间相较 Wi-Fi 对照方案减少两倍以上\parencite{retroar}。这些结果表明光的方向性能够保留用户的空间选择信息，并减少特定任务中的识别/连接开销，但不可脱离原实验条件外推。

RetroAR 的代价是用户必须拿起、移动并维持手机对准。手机光束可能与人想看的位置接近，但仍由手部动作间接表达关注。这一局限提出下一阶段问题：\textbf{如果方向可以选中对象，能否不用手持设备来表达这个方向？}

\section{EchoSight：从手持指向到佩戴式观察（已完成）}
\subsection{注意力线索与系统设计}
EchoSight 将交互入口从手机迁移至智能眼镜\parencite{echosight}。眼镜端 SightReader 与目标端 EchoTag 构成双向光学链路；双元件近红外光学设计及对齐检测利用空间方向抑制非目标响应，使用户在观察目标时获取信息并发送控制数据。系统延续了 RetroAR 的“方向性可选择目标”思想，但把指向动作尽量嵌入日常佩戴和观察，而不再要求每次举起手机。

局部光学链路把对象身份和信息交换放在同一个现场通道中，减少预注册与额外网络连接需求；双向能力又使用户在识别后能直接操控实体对象，而不止于读取标签。这是从“选择线索”向“原位交互”深入的关键。不过，EchoSight 的方向线索主要来自眼镜前端光路与用户观察方向的近似关系，\textbf{并非实测眼动追踪}。

\subsection{已有证据与通向下一阶段的局限}
CHI 2025 论文报告最远约 $5\,\mathrm{m}$ 的房间尺度通信和最大约 $120^\circ$ 的观察视角；其 12 名参与者的用户研究，在所设任务及对照条件下，显示相对二维码扫描、语音控制至少五倍的交互延迟降低，并报告更低的主观负担\parencite{echosight}。这些证据支持佩戴式、原位双向交互的效率与体验优势，但不能据此证明眼镜始终准确知道用户眼球注视点。

进一步的局限有三点：眼镜机身朝向和眼球实际视线可能偏离，固定光路仍可能迫使人调整姿态；用户看见某物可能只是观察而非操作；遮挡或转头可中断全光会话。由此引出第三阶段的核心问题：\textbf{能否让系统跟随真实视线，并可靠地区分“看见”与“想操作”？}

\section{Omnisight：从观察线索到视线对齐与意图判别（拟开展）}
Omnisight 是拟开展的毕业研究，以下模块均作为待实现、待验证的设计，不作为已经测得的系统结论。其总体思路是把先前依赖设备方向的目标选择推进为实测眼动驱动的对齐，再把目标定位与意图确认分开，并用不同链路分别承担选择和持续会话。

\subsection{实测眼动驱动的目标对齐}
首先，拟在眼镜端获取眼球视线并驱动可调光学发射单元，使光束方向尽可能跟随关注目标。系统需要标定眼动坐标、眼镜/头部姿态和光束坐标之间的变换，并在不同目标距离、相邻目标间距、头部运动和遮挡下测量对齐误差。将与固定光路方案比较目标选错率、对齐时延、功耗和佩戴负担，判定主动调节是否值得增加硬件复杂度。如果眼动置信度低，系统应暂缓执行并请求确认，而非把不确定目标当作正确目标。

\subsection{眼动与脑电的意图判别}
其次，拟在“看向但不操作”“主动选择目标”“执行确认”等可重复任务中，同步采集眼动与凝胶电极脑电，构建与实际交互动作相联系的标签。眼动用于定位候选对象；脑电仅被视为可能有增量价值的判别证据，需要处理眼动和肌电伪迹、个体差异与跨场景迁移。研究将比较仅眼动、仅脑电、融合方法以及眼动加显式确认的基线，报告误触发率、漏检率、精确率、召回率和决策时延，而不是只报告离线分类准确率。

脑电是否成为可用交互通道，取决于其增益是否覆盖采集、佩戴和计算成本。如果在控制眼动信息后，EEG 对误触发的改善不稳定或不足，系统将采用眼动加显式确认的备选路线，并把这一边界作为研究结论。

\subsection{定向选择与持续会话的分层链路}
最后，拟用下行定向光完成目标发现与身份确认，再经授权把该物理目标绑定到蓝牙设备身份；蓝牙承担上行反馈和后续持续会话。光负责“是哪一个”，蓝牙负责“如何保持交互”，两者之间需设计目标标识、状态机、超时、切换和重绑定机制，避免选中物体与实际控制设备不一致。实验应分别测试静止、转头、遮挡和移动场景中的绑定错误率、会话恢复时间、端到端延迟及能耗。混合设计的价值须以这些指标证明，不能假设蓝牙天然免配对或零开销。

目前的工作描述以研究问题、总体架构及关键器件方向为主；眼镜原型、脑电同步、融合模型和混合链路协议仍待实现或系统评估。最终交付应包括可复现实验原型、清晰的对照与消融、完整任务用户研究，以及对收益和成本的实证边界。

\section{综合比较：注意力利用如何逐级深入}
表\ref{tab:progression}把三项工作放在同一交互链上，区分已验证能力与拟检验能力。第一条递进路径是\textbf{注意力线索越来越贴近人}：手机方向、眼镜方向、实测视线与任务相关生理证据。第二条路径是\textbf{系统利用越来越接近行动决策}：先选中物体，再支持原位双向交互，最后检验是否真的应执行控制。第三条路径是\textbf{人对设备的适应负担逐步减少}：从举机对准，经过眼镜近似对齐，走向光束主动追随视线。

\begin{table}[htbp]
\centering\zihao{5}
\caption{三项工作对注意力利用的递进关系}
\label{tab:progression}
\begin{tabular}{p{2.0cm}p{2.8cm}p{4.7cm}p{3.7cm}}
\hline
系统 & 主要线索 & 已解决或拟解决的环节 & 尚需面对的边界 \\
\hline
RetroAR & 手持设备指向 & 已验证目标绑定与局部读取 & 持机、手动对准 \\
EchoSight & 眼镜光路近似观察 & 已验证原位双向交互 & 视线偏离、意图二义性 \\
Omnisight（拟） & 眼动与脑电 & 拟验证主动对齐、意图判别 & 增益、误触发、佩戴与链路开销 \\
\hline
\end{tabular}
\end{table}

统一框架并不意味着对所有场景都自动有效。只有空间线索能够可靠约束候选目标、实体与设备身份可验证绑定、意图确认的收益超过其成本时，注意力驱动交互才能真正改善体验。因此，最终评价必须比较完整任务中的时间、正确性与负担，而不是把三篇论文中不同实验条件下的单项数值简单排列。
